\chapter{संकेत और तरंगें}\label{ch:g10:signals-and-waves}

आपको आज तक मिला हर संदेश किसी संकेत के रूप में ही आया है: कोई
आवाज़ आपके कर्णपटह पर दबकर, और यह पन्ना रेडियो की एक झिलमिल बनकर। इस
विविधता के पीछे एक ही विचार बैठा है --- समय के साथ बदलती कोई
\omterm{def:g10:orders-of-magnitude:unit}{भौतिक राशि} --- और एक ही औज़ार-पेटी: \omterm{def:g10:signals-and-waves:period}{आवर्तकाल},
आवृत्ति, आयाम, और किसी देरी को नापने की कला। इसी से एक
जहाज़ समुद्र का तल नाप लेता है और एक चिकित्सक धड़कनें गिन लेता है।

\section{संकेत जानकारी ढोते हैं}

\begin{definition}[संकेत]\label{def:g10:signals-and-waves:signal}
\emph{संकेत}\index{संकेत} वह \omterm{def:g10:orders-of-magnitude:unit}{भौतिक राशि} है जो जानकारी ढोने के लिए
समय के साथ बदलती है: आपके कर्णपटह पर वायु का दाब, माइक्रोफ़ोन के
सिरों पर विभव, या किसी \omterm{def:g10:refraction:fiber}{प्रकाशिक तंतु} में प्रकाश की चमक। किसी संकेत
का अध्ययन करने के लिए उसे समय के फलन के रूप में दर्ज कीजिए और जानकारी
वक्र से ही पढ़ लीजिए।
\end{definition}

किसी फ़ोन कॉल में वही जानकारी बार-बार वाहक बदलती है: पहले वायु में
दाब, फिर माइक्रोफ़ोन में विभव, फिर एंटीना पर रेडियो --- और दूसरे
सिरे पर उलटे क्रम में। इन्हें एक ही शब्दावली सँभाल लेती है; यह अध्याय
वही गढ़ता है।

\section{आवर्ती संकेत}

\begin{definition}[आवर्ती संकेत, आवर्तकाल]\label{def:g10:signals-and-waves:period}
कोई संकेत \emph{आवर्ती}\index{आवर्ती संकेत} तब होता है जब वह
नियमित अंतरालों पर हूबहू दोहराया जाए। \emph{आवर्तकाल}\index{आवर्तकाल}
$T$ एक पूरे चक्र की अवधि है, सेकंडों में।
\end{definition}

\begin{definition}[आवृत्ति]\label{def:g10:signals-and-waves:frequency}
किसी \omterm{def:g10:signals-and-waves:period}{आवर्ती संकेत} की \emph{आवृत्ति}\index{आवृत्ति} प्रति \omterm{def:g10:orders-of-magnitude:unit}{सेकंड}
चक्रों की संख्या है:
\[
f = \frac{1}{T} .
\]
इसका मात्रक, यानी एक चक्र प्रति \omterm{def:g10:orders-of-magnitude:unit}{सेकंड},
\emph{हर्ट्ज़}\index{हर्ट्ज़} (\unit{Hz}) कहलाता है, और उसके सामान्य गुणज
\unit{kHz}, \unit{MHz}, \unit{GHz} हैं।
\end{definition}

\begin{definition}[आयाम]\label{def:g10:signals-and-waves:amplitude}
किसी \omterm{def:g10:signals-and-waves:period}{आवर्ती संकेत} का \emph{आयाम}\index{आयाम} विश्राम-मान से उसका
अधिकतम विचलन है। शून्य के इर्द-गिर्द सममित रूप से दोलन करते विभव के लिए
आयाम शिखर-मान होता है; शिखर से गर्त तक की नाप आयाम का \emph{दुगुना}
(शिखर-से-शिखर मान) देती है।
\end{definition}

\begin{example}[विद्युत आपूर्ति और संगीत का सुर]\label{ex:g10:signals-and-waves:mains}
घरेलू विद्युत आपूर्ति का विभव $f = \qty{50}{Hz}$ पर दोलन करता है: उसका
\omterm{def:g10:signals-and-waves:period}{आवर्तकाल} $T = 1/f = \qty{0.020}{s} = \qty{20}{ms}$ है। किसी वाद्यवृंद का सुर A \qty{440}{Hz} पर एक
दाब संकेत है: उसका एक चक्र $T = 1/440 = \qty{2.3}{ms}$ चलता है। आवृत्ति
जितनी अधिक, चक्र उतना छोटा: $f$ और $T$ एक-दूसरे के व्युत्क्रम हैं।
\end{example}

\begin{omfigure}
\begin{tikzpicture}[font=\small]
  \draw[-Stealth] (-0.3,0) -- (7.7,0) node[below] {$t$};
  \draw[-Stealth] (0,-1.7) -- (0,2.0) node[left] {संकेत};
  \draw[omDef, thick, smooth, samples=120, domain=0:7.3]
    plot (\x, {1.3*sin(120*\x)});
  \draw[dashed, thin] (0.75,1.3) -- (0.75,1.6) (3.75,1.3) -- (3.75,1.6);
  \draw[omThm, Stealth-Stealth] (0.75,1.6) -- (3.75,1.6)
    node[midway, above] {आवर्तकाल $T$};
  \draw[omProp, Stealth-Stealth] (5.25,0) -- (5.25,-1.3)
    node[midway, right] {आयाम};
\end{tikzpicture}

{\small एक \omterm{def:g10:signals-and-waves:period}{आवर्ती संकेत}: पैटर्न हर $T$ \omterm{def:g10:orders-of-magnitude:unit}{सेकंड} बाद दोहराया जाता
है, और आयाम विश्राम-मान से अधिकतम विचलन है।}
\end{omfigure}

\section{दोलनलेख पढ़ना}

\begin{definition}[दोलनदर्शी]\label{def:g10:signals-and-waves:oscilloscope}
\emph{दोलनदर्शी}\index{दोलनदर्शी} किसी विभव का समय के सामने आलेख खींचता
है; उसके जालीदार परदे पर बना वक्र
\emph{दोलनलेख}%\index{दोलनलेख} कहलाता है। दो सेटिंग परदे के खानों
को भौतिक मानों में बदल देती हैं: \emph{कालाधार}\index{कालाधार} (प्रति
क्षैतिज खाना \omterm{def:g10:orders-of-magnitude:unit}{सेकंड}) और \omterm{def:g10:universal-gravitation:weight}{ऊर्ध्वाधर} \emph{लब्धि}\index{लब्धि} (प्रति
\omterm{def:g10:universal-gravitation:weight}{ऊर्ध्वाधर} खाना वोल्ट)।
\end{definition}

\begin{method}[परदे से संख्याओं तक]\label{met:g10:signals-and-waves:oscillogram}
\begin{enumerate}
  \item एक पूरे चक्र में समाए क्षैतिज खाने गिनिए (शिखर से शिखर तक गिनना
        सबसे सुरक्षित है); उन्हें \omterm{def:g10:signals-and-waves:oscilloscope}{कालाधार} से गुणा करके $T$
        निकालिए, फिर $f = 1/T$।
  \item मध्य रेखा से किसी शिखर तक \omterm{def:g10:universal-gravitation:weight}{ऊर्ध्वाधर} खाने गिनिए और उन्हें
        \omterm{def:g10:signals-and-waves:oscilloscope}{लब्धि} से गुणा करके आयाम निकालिए (शिखर से गर्त तक
        गिनने पर शिखर-से-शिखर मान मिलता है: उसे आधा कर लीजिए)।
\end{enumerate}
\end{method}

\begin{example}[एक पूरा पाठ]\label{ex:g10:signals-and-waves:reading}
नीचे दिए \omterm{def:g10:signals-and-waves:oscilloscope}{दोलनलेख} में \omterm{def:g10:signals-and-waves:oscilloscope}{कालाधार} \qty{5}{ms} प्रति खाना है और
\omterm{def:g10:signals-and-waves:oscilloscope}{लब्धि} \qty{2}{V} प्रति खाना। एक चक्र $4.0$ खानों में फैला है:
$T = 4.0 \times \qty{5}{ms} = \qty{20}{ms}$ और $f = 1/\qty{0.020}{s} = \qty{50}{Hz}$ --- यानी फिर वही घरेलू आपूर्ति। एक शिखर अक्ष से $3.0$
खाने ऊपर बैठा है: आयाम $3.0 \times \qty{2}{V} = \qty{6.0}{V}$ है।
\end{example}

\begin{omfigure}
\begin{tikzpicture}[x=0.62cm, y=0.62cm, font=\small]
  \draw[gray!35, thin] (0,0) grid (10,8);
  \draw[gray!80] (0,4) -- (10,4);
  \draw[gray!80] (5,0) -- (5,8);
  \draw[omDef, thick, smooth, samples=120, domain=0:10]
    plot (\x, {4 + 3*sin(90*\x)});
  \draw[dashed, thin] (1,7) -- (1,7.6) (5,7) -- (5,7.6);
  \draw[omThm, Stealth-Stealth] (1,7.6) -- (5,7.6)
    node[midway, above] {$4.0$ खाने};
  \draw[omProp, Stealth-Stealth] (7,4) -- (7,1)
    node[midway, right] {$3.0$ खाने};
\end{tikzpicture}

{\small एक \omterm{def:g10:signals-and-waves:oscilloscope}{दोलनलेख}, \omterm{def:g10:signals-and-waves:oscilloscope}{कालाधार} \qty{5}{ms}/खाना और \omterm{def:g10:signals-and-waves:oscilloscope}{लब्धि}
\qty{2}{V}/खाना: एक चक्र $4.0$ खानों में फैला है ($T = \qty{20}{ms}$, $f = \qty{50}{Hz}$) और शिखर
$3.0$ खाने ऊपर है (आयाम \qty{6.0}{V})।}
\end{omfigure}

\section{ध्वनि और विद्युतचुंबकीय संकेत}

\begin{definition}[ध्वनि]\label{def:g10:signals-and-waves:sound}
\emph{ध्वनि}\index{ध्वनि} एक दाब संकेत है: कोई कंपन करती
वस्तु --- तार, झिल्ली, स्वर-रज्जु --- वायु को लयबद्ध ढंग से धकेलती है,
और संपीड़न बाहर की ओर लगभग \qty{340}{m/s} से चलते हैं। माइक्रोफ़ोन आते हुए
दाब को विभव में बदल देता है; कर्णपटह उसे तंत्रिका-स्पंदों में।
ध्वनि को माध्यम चाहिए: निर्वात के जार में बजती घंटी वायु निकलते-निकलते
चुप हो जाती है।
\end{definition}

\begin{definition}[तरंग]\label{def:g10:signals-and-waves:wave}
चलता हुआ कोई विक्षोभ --- \omterm{def:g10:signals-and-waves:sound}{ध्वनि} का संपीड़न-पैटर्न, तालाब की लहर,
या रेडियो का कोई स्पंद --- \emph{तरंग}\index{तरंग} कहलाता है। माध्यम, जब
हो, अपनी जगह ही रहता है: लहर तालाब पार करती है जबकि पानी का हर टुकड़ा
वहीं ऊपर-नीचे डोलता रहता है। तरंगों की ज्यामिति --- \omterm{def:g10:light-spectra:wavelength}{तरंगदैर्घ्य},
व्यतिकरण, विवर्तन --- आगे के किसी अध्याय में उठाई जाएगी; इस अध्याय को
केवल उनकी चालें चाहिए।
\end{definition}

\begin{definition}[पराश्रव्य और अवश्रव्य]\label{def:g10:signals-and-waves:ultrasound}
मनुष्य का कान लगभग \qty{20}{Hz} से \qty{20}{kHz} के बीच की \omterm{def:g10:signals-and-waves:sound}{ध्वनि} सुनता है
(ऊपरी सीमा उम्र के साथ गिरती जाती है)। \qty{20}{kHz} से ऊपर की \omterm{def:g10:signals-and-waves:sound}{ध्वनि}
\emph{पराश्रव्य}\index{पराश्रव्य} है --- चमगादड़ \qty{50}{kHz} पर शिकार करते
हैं, और चिकित्सा के स्कैनर कुछ मेगाहर्ट्ज़ पर शरीर के चित्र लेते हैं।
\qty{20}{Hz} से नीचे की \omterm{def:g10:signals-and-waves:sound}{ध्वनि} \emph{अवश्रव्य}\index{अवश्रव्य} है, जिसे
हाथी और भूकंपमापी महसूस कर लेते हैं।
\end{definition}

\begin{proposition}[विद्युतचुंबकीय संकेत]\label{prop:g10:signals-and-waves:em}
रेडियो, वाई-फ़ाई और प्रकाश \emph{विद्युतचुंबकीय
संकेत}\index{विद्युतचुंबकीय संकेत} हैं: एक ही कुल, सभी निर्वात से होकर
चलते हैं --- किसी माध्यम की ज़रूरत नहीं --- और सभी एक ही चाल से, यानी
\cref{ch:g10:refraction} की प्रकाश-चाल से:
\[
c = \qty{3.00e8}{m/s} \quad \text{(वायु में भी लगभग यही)।}
\]
\end{proposition}
\admitted

\begin{remark}\label{rem:g10:signals-and-waves:honest}
यहाँ यह एक प्रायोगिक तथ्य है --- धूप खाली अंतरिक्ष के $150$ करोड़
किलोमीटर पार कर आती है; इसकी ईमानदार व्युत्पत्ति स्नातक वर्ष 2 के खंड
में विद्युत और चुंबकत्व के नियमों से की गई है। यह भेद याद रखिए:
\omterm{def:g10:signals-and-waves:sound}{ध्वनि} माध्यम \emph{का} कंपन है और उसके बिना मर जाती है; जबकि
\omterm{prop:g10:signals-and-waves:em}{विद्युतचुंबकीय संकेत} अपना क्षेत्र अपने साथ निर्वात के पार ले
जाता है, और वह भी दस लाख गुना तेज़।
\end{remark}

\begin{omfigure}
\begin{tikzpicture}[x=1.3cm, font=\small]
  \fill[omDef!15] (1.3,-0.14) rectangle (4.3,0.14);
  \draw[-Stealth] (-0.3,0) -- (7.7,0);
  \foreach \x/\l in {0/{$1$}, 1/{$10$}, 2/{$10^2$}, 3/{$10^3$},
      4/{$10^4$}, 5/{$10^5$}, 6/{$10^6$}, 7/{$10^7$}}
    \draw (\x,0.07) -- (\x,-0.07) node[below] {\l};
  \node[below] at (7.55,-0.07) {Hz};
  \node[gray] at (0.4,-0.75) {अवश्रव्य};
  \node[omDef] at (2.8,-0.75) {श्रव्य};
  \node[gray] at (5.8,-0.75) {पराश्रव्य};
  \draw[omThm] (1.7,0.14) -- (1.7,0.4) node[above] {आपूर्ति};
  \draw[omThm] (2.64,0.14) -- (2.64,0.95) node[above] {सुर A};
  \draw[omThm] (4.7,0.14) -- (4.7,0.4) node[above] {चमगादड़};
  \draw[omThm] (6.7,0.14) -- (6.7,0.95) node[above] {चिकित्सा स्कैन};
\end{tikzpicture}

{\small लघुगणकीय पैमाने पर \omterm{def:g10:signals-and-waves:sound}{ध्वनि} की आवृत्तियाँ (हर क़दम $\times 10$)।
कान \qty{20}{Hz} से \qty{20}{kHz} तक सुनता है; नीचे \omterm{def:g10:signals-and-waves:ultrasound}{अवश्रव्य} है और ऊपर
\omterm{def:g10:signals-and-waves:ultrasound}{पराश्रव्य}।}
\end{omfigure}

\section{प्रतिध्वनि: देरी से दूरी}

$v$ चाल से चलता कोई संकेत $t$ समय में $d = v\,t$ तय करता
है: इसलिए हर देरी भेस बदली हुई दूरी है। इससे भी अच्छा: किसी रुकावट पर एक
छोटा स्पंद भेजिए और उसकी \emph{प्रतिध्वनि}\index{प्रतिध्वनि} का समय
नापिए --- आने-जाने में वह $2d$ तय करता है।

\begin{method}[प्रतिध्वनि से परास नापना]\label{met:g10:signals-and-waves:echo}
\begin{enumerate}
  \item एक छोटा स्पंद भेजिए; \omterm{rem:g10:signals-and-waves:honest}{प्रतिध्वनि} के आने-जाने का समय
        $t$ नापिए।
  \item \emph{जिस माध्यम को पार किया गया है उसमें} संकेत की चाल
        $v$ लीजिए।
  \item स्पंद ने $2d = v\,t$ तय किया, इसलिए $d = \dfrac{v\,t}{2}$।
\end{enumerate}
पानी में \omterm{def:g10:signals-and-waves:sound}{ध्वनि} के स्पंद \emph{सोनार}\index{सोनार} बनाते हैं
($v \approx \qty{1500}{m/s}$); वायु में रेडियो स्पंद \emph{रडार}\index{रडार} ($v = c$); और
शरीर में \omterm{def:g10:signals-and-waves:ultrasound}{पराश्रव्य} स्पंद चिकित्सा का स्कैनर (मुलायम ऊतक में
$v \approx \qty{1540}{m/s}$)।
\end{method}

\begin{example}[सोनार]\label{ex:g10:signals-and-waves:sonar}
किसी जहाज़ का \omterm{met:g10:signals-and-waves:echo}{सोनार} एक पिंग भेजता है और समुद्र-तल की
\omterm{rem:g10:signals-and-waves:honest}{प्रतिध्वनि} $t = \qty{0.80}{s}$ बाद लौटती है: गहराई $d = \frac{1500 \times 0.80}{2} = \qty{600}{m}$ है। गुणक $2$
भूल जाइए और महासागर दुगुना हो जाएगा।
\end{example}

\begin{omfigure}
\begin{tikzpicture}[font=\small]
  \fill[blue!8] (-3.5,0) -- (3.5,0) -- (3.5,-2.95) -- (2.6,-3.1)
    -- (1.4,-2.9) -- (0.2,-3.05) -- (-1.2,-2.85) -- (-2.3,-3.1)
    -- (-3.5,-2.95) -- cycle;
  \draw[thick, blue!60] (-3.5,0) -- (3.5,0);
  \draw[thick] (-3.5,-2.95) -- (-2.3,-3.1) -- (-1.2,-2.85)
    -- (0.2,-3.05) -- (1.4,-2.9) -- (2.6,-3.1) -- (3.5,-2.95);
  \fill[pattern=north east lines] (-3.5,-2.95) -- (-2.3,-3.1)
    -- (-1.2,-2.85) -- (0.2,-3.05) -- (1.4,-2.9) -- (2.6,-3.1)
    -- (3.5,-2.95) -- (3.5,-3.4) -- (-3.5,-3.4) -- cycle;
  \fill[gray!70] (-0.9,0.3) -- (0.9,0.3) -- (0.55,-0.25)
    -- (-0.55,-0.25) -- cycle;
  \fill[gray!50] (-0.25,0.3) rectangle (0.25,0.62);
  \draw[-Stealth, omDef, thick] (-0.35,-0.35) -- (-0.35,-2.75)
    node[midway, left] {पिंग, $t/2$};
  \draw[-Stealth, omThm, thick] (0.35,-2.75) -- (0.35,-0.35)
    node[midway, right] {प्रतिध्वनि, $t/2$};
  \draw[Stealth-Stealth] (2.9,-0.05) -- (2.9,-2.9)
    node[midway, right] {$d = \dfrac{v\,t}{2}$};
\end{tikzpicture}

{\small \omterm{met:g10:signals-and-waves:echo}{सोनार} का पिंग नीचे जाकर लौटता है: नापी गई देरी $t$
में $2d$ तय होता है, इसलिए समुद्री पानी में $v \approx \qty{1500}{m/s}$ लेकर गहराई $d = v\,t/2$
होती है।}
\end{omfigure}

\begin{example}[रडार]\label{ex:g10:signals-and-waves:radar}
हवाई अड्डे का \omterm{met:g10:signals-and-waves:echo}{रडार} किसी विमान की \omterm{rem:g10:signals-and-waves:honest}{प्रतिध्वनि} स्पंद के
$\qty{0.30}{ms}$ बाद पाता है: $d = \frac{\qty{3.00e8}{m/s} \times \qty{3.0e-4}{s}}{2} = \qty{45}{km}$। $c$ इतना बड़ा है कि \omterm{met:g10:signals-and-waves:echo}{रडार}
माइक्रोसेकंड की घड़ियों पर जीता है: प्रकाश का समय नापना बारीक काम है।
\end{example}

\begin{example}[संकेत के रूप में हृदय]\label{ex:g10:signals-and-waves:ecg}
हर धड़कन छाती के आरपार एक छोटा विभव-स्पंद भेजती है; इलेक्ट्रोड उसे
\emph{विद्युत-हृद्लेख}\index{विद्युत-हृद्लेख} (ECG) के रूप में दर्ज कर
लेते हैं, जो लगभग एक \omterm{def:g10:signals-and-waves:period}{आवर्ती संकेत} है। यदि ऊँची चोटियाँ $T = \qty{0.75}{s}$
की दूरी पर हों तो हृदय $f = 1/0.75 = \qty{1.33}{Hz}$ पर, यानी $1.33 \times 60 = 80$ धड़कन प्रति मिनट पर
धड़क रहा है: किसी लेख से $T$ पढ़ लेना ही वह तरीक़ा है जिससे
हृदय-रोग विशेषज्ञ आपकी नाड़ी लेता है।
\end{example}

\section{अभ्यास}

\begin{exercise}[$\star$]\label{exo:g10:signals-and-waves:1}
निम्नलिखित आवर्तकालों वाले \omterm{def:g10:signals-and-waves:period}{आवर्ती} संकेतों की आवृत्ति
निकालिए: (क) $T = \qty{20}{ms}$; (ख) $T = \qty{4.0}{ms}$; (ग) $T = \qty{50}{\micro s}$।
\end{exercise}

\begin{exercise}[$\star$]\label{exo:g10:signals-and-waves:2}
इनका \omterm{def:g10:signals-and-waves:period}{आवर्तकाल} निकालिए: (क) \qty{50}{Hz} की घरेलू आपूर्ति; (ख)
\qty{440}{Hz} का सुर A; (ग) \qty{2.4}{GHz} का वाई-फ़ाई संकेत।
\end{exercise}

\begin{exercise}[$\star$]\label{exo:g10:signals-and-waves:3}
इन्हें \omterm{def:g10:signals-and-waves:ultrasound}{अवश्रव्य}, श्रव्य \omterm{def:g10:signals-and-waves:sound}{ध्वनि} या \omterm{def:g10:signals-and-waves:ultrasound}{पराश्रव्य} में
बाँटिए: \qty{12}{Hz}; \qty{440}{Hz}; \qty{18}{kHz}; \qty{40}{kHz} (कार का पार्किंग संवेदक);
\qty{5}{MHz} (चिकित्सा की जाँच-छड़)।
\end{exercise}

\begin{exercise}[$\star$]\label{exo:g10:signals-and-waves:4}
प्रति खाना \qty{2}{ms} और प्रति खाना \qty{0.5}{V} पर सेट किए \omterm{def:g10:signals-and-waves:oscilloscope}{दोलनदर्शी}
में एक चक्र $5.0$ खानों में फैला है और शिखर अक्ष से $4.0$ खाने ऊपर
बैठा है। $T$, $f$ और आयाम ज्ञात कीजिए।
\end{exercise}

\begin{exercise}[$\star$]\label{exo:g10:signals-and-waves:5}
आप बिजली की चमक देखते हैं और \qty{6.0}{s} बाद गरज सुनते हैं। बिजली कितनी दूर
गिरी ($v_{\text{ध्वनि}} = \qty{340}{m/s}$)? प्रकाश का यात्रा-समय क्यों छोड़ा जा सकता है?
\end{exercise}

\begin{exercise}[$\star\star$]\label{exo:g10:signals-and-waves:6}
समुद्र-तल से \omterm{met:g10:signals-and-waves:echo}{सोनार} की \omterm{rem:g10:signals-and-waves:honest}{प्रतिध्वनि} \qty{0.60}{s} बाद लौटती है
(समुद्री पानी में $v = \qty{1500}{m/s}$)। गुणनफल $v\,t$ को आधा क्यों करना पड़ता है?
गहराई निकालिए। \qty{1200}{m} गहरी किसी खाई पर \omterm{rem:g10:signals-and-waves:honest}{प्रतिध्वनि} को कितना
समय लगेगा?
\end{exercise}

\begin{exercise}[$\star\star$]\label{exo:g10:signals-and-waves:7}
हवाई अड्डे का \omterm{met:g10:signals-and-waves:echo}{रडार} किसी विमान की \omterm{rem:g10:signals-and-waves:honest}{प्रतिध्वनि} \qty{0.24}{ms} बाद
पाता है और ठीक \qty{1.0}{s} बाद एक दूसरी \omterm{rem:g10:signals-and-waves:honest}{प्रतिध्वनि} \qty{0.238}{ms} बाद।
दोनों दूरियाँ निकालिए, फिर विमान की चाल। क्या वह पास आ रहा है?
\end{exercise}

\begin{exercise}[$\star\star$]\label{exo:g10:signals-and-waves:8}
कोई ECG \qty{25}{mm/s} पर छापा जाता है; ऊँची चोटियाँ \qty{20}{mm} की दूरी पर हैं।
\omterm{def:g10:signals-and-waves:period}{आवर्तकाल} और हृदय-गति (धड़कन प्रति मिनट) ज्ञात कीजिए।
क्षिप्रहृदयता का अर्थ है $100$ धड़कन प्रति मिनट से अधिक: इस काग़ज़ पर
चोटियों की दूरी किससे कम होने पर वह दिखेगी?
\end{exercise}

\begin{exercise}[$\star\star$]\label{exo:g10:signals-and-waves:9}
आपको \qty{3.0}{V} आयाम वाला \qty{200}{Hz} का एक संकेत ऐसे परदे पर
दिखाना है जिसमें $10$ क्षैतिज और $8$ \omterm{def:g10:universal-gravitation:weight}{ऊर्ध्वाधर} खाने हैं (मध्य
रेखा से ऊपर $4$)। ऐसा \omterm{def:g10:signals-and-waves:oscilloscope}{कालाधार} चुनिए जो लगभग दो पूरे चक्र
दिखाए, और ऐसी \omterm{def:g10:signals-and-waves:oscilloscope}{लब्धि} जो बिना कटे परदे का अधिकांश हिस्सा काम में
ले।
\end{exercise}

\begin{exercise}[$\star\star$]\label{exo:g10:signals-and-waves:10}
चिकित्सा की एक जाँच-छड़ पेट में \omterm{def:g10:signals-and-waves:ultrasound}{पराश्रव्य} का एक स्पंद भेजती है
(मुलायम ऊतक में $v = \qty{1540}{m/s}$) और \qty{40}{\micro s} (किसी अंग की अगली दीवार) तथा \qty{60}{\micro s}
(पिछली दीवार) बाद प्रतिध्वनियाँ सुनती है। हर दीवार की गहराई और
अंग की मोटाई ज्ञात कीजिए।
\end{exercise}

\begin{exercise}[$\star\star$]\label{exo:g10:signals-and-waves:11}
किसी संगीत-कार्यक्रम का सीधा प्रसारण हो रहा है। ढोल की थाप पहले कौन
सुनेगा: मंच से \qty{30}{m} दूर बैठा श्रोता, या \qty{3000}{km} दूर रेडियो सुनता
श्रोता? और कितना पहले?
\end{exercise}

\begin{exercise}[$\star\star\star$]\label{exo:g10:signals-and-waves:12}
एक चमगादड़ \qty{3.0}{ms} तक चलने वाले \omterm{def:g10:signals-and-waves:ultrasound}{पराश्रव्य} स्पंद छोड़ता है और
जब तक छोड़ता रहता है तब तक कोई \omterm{rem:g10:signals-and-waves:honest}{प्रतिध्वनि} सुन नहीं सकता। किस
दूरी से पास होने पर पतंगा पकड़ में नहीं आएगा? \qty{1.7}{m} दूर बैठे पतंगे की
\omterm{rem:g10:signals-and-waves:honest}{प्रतिध्वनि} की देरी कितनी है? अंतिम धावे के समय चमगादड़ को अपने
स्पंद छोटे क्यों करने पड़ते हैं?
\end{exercise}

\begin{exercise}[$\star\star\star$]\label{exo:g10:signals-and-waves:13}
आप किसी चट्टान की ओर मुँह करके ताली बजाते हैं और \qty{1.5}{s} बाद
\omterm{rem:g10:signals-and-waves:honest}{प्रतिध्वनि} सुनते हैं। चट्टान कितनी दूर है? आप सीधे उसकी ओर
\qty{100}{m} चल देते हैं: नई देरी? लगभग \qty{0.1}{s} से कम पर कान ताली और
\omterm{rem:g10:signals-and-waves:honest}{प्रतिध्वनि} को अलग नहीं कर पाता: किस दूरी के भीतर
\omterm{rem:g10:signals-and-waves:honest}{प्रतिध्वनि} ताली में ही घुल जाएगी?
\end{exercise}

\begin{exercise}[$\star\star\star$]\label{exo:g10:signals-and-waves:14}
प्रति खाना \qty{5}{ms} पर किसी संकेत का एक चक्र $10$ खानों वाले
परदे के $3.0$ खानों में फैला है।
\begin{enumerate}
  \item $T$ और $f$ ज्ञात कीजिए।
  \item \omterm{def:g10:signals-and-waves:oscilloscope}{कालाधार} बदलकर प्रति खाना \qty{2}{ms} कर दिया जाता है: अब एक
        चक्र कितने खानों में फैलेगा?
  \item हर सेटिंग पर परदे में कितने \emph{पूरे} चक्र समाते हैं?
\end{enumerate}
\end{exercise}

\begin{exercise}[$\star\star\star$]\label{exo:g10:signals-and-waves:15}
कोई ``पिंग'' \qty{1200}{km} दूर बैठे किसी सर्वर तक इंटरनेट के एक पैकेट के
आने-जाने का समय नापता है। सबसे कम संभव आने-जाने का समय क्या है
($c$ पर चलते संकेत)? \omterm{def:g10:refraction:fiber}{प्रकाशिक तंतु} में प्रकाश लगभग
\qty{2.0e8}{m/s} से चलता है: फिर से निकालिए। नापा गया पिंग \qty{40}{ms} है: वह आपके
उत्तरों से अधिक क्यों है, इसके दो कारण दीजिए।
\end{exercise}

\section{समस्या: सोनार, तूफ़ान और हृद्लेख}

\begin{problem}[{सप्ताहांत समस्या --- दुनिया के संकेत पढ़ना: एक ही नियम,
$d = v\,t$, जो तूफ़ान की दूरी नापता है, समुद्र-तल का नक्शा बनाता है, धड़कनें
गिनता है और GPS का भेद भी खोल देता है}]
\label{pb:g10:signals-and-waves:1}
मछली पकड़ने वाला एक जहाज़ तूफ़ानी रात में काम कर रहा है: क्षितिज पर
बिजली, तल पर झाड़ू फेरता \omterm{met:g10:signals-and-waves:echo}{सोनार}, चिकित्सक की काग़ज़-पट्टी पर
कप्तान का हृदय, और उपग्रहों को सुनता एक GPS ग्राही। चार उपकरण,
एक नियम। वायु में $v_{\text{ध्वनि}} = \qty{340}{m/s}$, समुद्री पानी में \qty{1500}{m/s}, और $c = \qty{3.00e8}{m/s}$ लीजिए।

\medskip
\noindent\textbf{भाग I --- तूफ़ान।}
\begin{enumerate}
  \item एक चमक; गरज \qty{9.0}{s} बाद आती है। वह कितनी दूर है?
  \item उस दूरी पर प्रकाश का यात्रा-समय निकालिए, और चमक को तात्क्षणिक
        मानना उचित ठहराइए।
  \item नाविक चमक और गरज के बीच के \omterm{def:g10:orders-of-magnitude:unit}{सेकंड} गिनकर तीन से भाग दे देते हैं और
        किलोमीटर पा जाते हैं। इस नियम को उचित ठहराइए।
  \item पाँच मिनट बाद एक चमक \qty{6.0}{s} देती है। अब कितनी दूर? तूफ़ान के
        पास आने की \omterm{def:g10:relative-motion:speed}{औसत चाल} \unit{m/s} और \unit{km/h} में निकालिए।
  \item गरज कड़कने के बजाय गड़गड़ाती है: बिजली की नाली किलोमीटरों लंबी
        होती है, इसलिए उसके हिस्से अलग-अलग दूरियों पर पड़ते हैं। यदि वह
        नाली जहाज़ से \qty{2.0}{km} से \qty{5.0}{km} तक फैली हो, तो गड़गड़ाहट कितनी
        देर चलेगी?
\end{enumerate}

\medskip
\noindent\textbf{भाग II --- \omterm{met:g10:signals-and-waves:echo}{सोनार}।}
\begin{enumerate}[resume]
  \item \omterm{met:g10:signals-and-waves:echo}{सोनार} पिंग भेजता है; समुद्र-तल \qty{0.90}{s} में जवाब देता है।
        गहराई?
  \item उथले हिस्से पर \omterm{rem:g10:signals-and-waves:honest}{प्रतिध्वनि} घटकर \qty{0.20}{s} रह जाती है।
        गहराई?
  \item एक ही पिंग \emph{दो} प्रतिध्वनियाँ लौटाता है, \qty{0.16}{s}
        और \qty{0.20}{s} पर। उनकी व्याख्या कीजिए; मछलियों का झुंड तल से कितना
        ऊपर तैर रहा है?
  \item \omterm{met:g10:signals-and-waves:echo}{सोनार} हर \qty{0.50}{s} पर पिंग भेजता है। बिना गड़बड़ाए वह
        अधिक से अधिक कितनी गहराई नाप सकता है, और उससे आगे क्या गड़बड़
        होती है?
  \item वायु में वही \qty{0.60}{s} की देरी किस दूरी का अर्थ रखती? सीख: किसी
        देरी को दूरी में बदलने से पहले आपको क्या पता होना चाहिए?
\end{enumerate}

\medskip
\noindent\textbf{भाग III --- हृद्लेख।}
\begin{enumerate}[resume]
  \item कप्तान के ECG की चोटियाँ \qty{0.86}{s} की दूरी पर हैं।
        आवृत्ति? धड़कन प्रति मिनट?
  \item पट्टी \qty{25}{mm/s} की दर से आगे बढ़ती है। चिकित्सक ने चोटियों के बीच
        कितनी दूरी नापी?
  \item जाल खींचने के बाद दूरी \qty{15}{mm} रह जाती है। नई हृद्गति?
  \item एक युवा नाविक, जो अभ्यस्त नाविक-खिलाड़ी है, विश्राम में $50$
        धड़कन प्रति मिनट पर रहता है। \omterm{def:g10:signals-and-waves:period}{आवर्तकाल}? पट्टी पर दूरी?
  \item क्या ECG कड़ाई से \omterm{def:g10:signals-and-waves:period}{आवर्ती} है? फिर चिकित्सक उस लेख से
        क्या पढ़ता है?
\end{enumerate}

\medskip
\noindent\textbf{भाग IV --- प्रकाश के सहारे घर।}
\begin{enumerate}[resume]
  \item जहाज़ \qty{50}{km} दूर बंदरगाह को रेडियो संदेश भेजता है। संदेश को
        कितना समय लगता है? \omterm{def:g10:signals-and-waves:sound}{ध्वनि} को कितना लगता?
  \item कोई GPS \omterm{def:g10:universal-gravitation:satellite}{उपग्रह} सिर के ऊपर लगभग \qty{20200}{km} पर
        कक्षा में है। उसके संकेत को जहाज़ तक पहुँचने
        में कितना समय लगता है?
  \item GPS समय को स्थिति में बदल देता है। \qty{1.0}{\micro s} की घड़ी-त्रुटि कितनी
        दूरी की त्रुटि पैदा करती है? \qty{3.0}{m} की सटीकता के लिए समय कितनी
        परिशुद्धि से नापना पड़ेगा?
  \item \omterm{def:g10:signals-and-waves:sound}{ध्वनि} पर आधारित GPS सैद्धांतिक रूप से भी क्यों नहीं बन सकता?
  \item समापन: आज रात काम आई तीनों चालें और चारों उपकरणों के पीछे का
        इकलौता नियम गिनाइए, फिर एक वाक्य में परास नापने की वह आदत बताइए।
\end{enumerate}
\end{problem}
